% Generated draft: requires mathematical and source review.
\documentclass[11pt]{article}
\usepackage[margin=25mm]{geometry}
\usepackage{fontspec}
\setmainfont{Latin Modern Roman}
\usepackage{amsmath,amssymb,mathtools,unicode-math}
\setmathfont{Latin Modern Math}
\usepackage{enumitem,needspace,xurl,hyperref}
\setlength{\parindent}{0pt}
\setlength{\parskip}{5pt}
\setlength{\emergencystretch}{3em}
\newcommand{\N}{\mathbb{N}}
\newcommand{\Z}{\mathbb{Z}}
\newcommand{\Q}{\mathbb{Q}}
\newcommand{\R}{\mathbb{R}}
\newcommand{\C}{\mathbb{C}}
\newcommand{\F}{\mathbb{F}}
\newcommand{\A}{\mathbb{A}}
\newcommand{\PP}{\mathbb{P}}
\newcommand{\E}{\mathbb{E}}
\newcommand{\eps}{\varepsilon}
\newcommand{\dd}{\,\mathrm{d}}
\begin{document}
\section*{UnsolvedMath ID: 10400059}
\textbf{Source problem code:} catalogue entry
\subsection{Definitions and mathematical statement}
A virtual knot is represented by a diagram with classical and virtual crossings, modulo classical Reidemeister moves and the virtual detour equivalence. Whether the knots here are closed or long must be determined from the source, since its target uses an interval \(I\). Write \(\overrightarrow{\mathcal A}(I)\) for exactly the source's directed-diagram algebra on \(I\); its diagram generators, relations and grading are not included in the supplied record.

M. Polyak's Problem 3.9 asks to construct a ``Kontsevich invariant'' of virtual knots, meaning a universal finite-type invariant, with values in \(\overrightarrow{\mathcal A}(I)\). To explain universality conditionally, suppose the source fixes a coefficient field \(F\), a virtual-knot class \(\mathcal V\), a finite-type filtration \(F^{r+1}(F\mathcal V)\), and a graded diagram target \(\mathcal D=\overrightarrow{\mathcal A}(I)\). A completed universal invariant would be an isotopy invariant
\[
 Z_{\rm virt}:\mathcal V\longrightarrow\widehat{\mathcal D}
\]
whose truncations encode all finite-type invariants: for each \(r\ge0\) and each \(F\)-valued invariant \(v\) vanishing on \(F^{r+1}(F\mathcal V)\), there is a linear functional \(\lambda_v\) on the degree-at-most-\(r\) target such that
\[
 v=\lambda_v\circ\operatorname{pr}_{\le r}\circ Z_{\rm virt}.
\]
A Kontsevich-type construction ordinarily also specifies its associated-graded symbol map. The required diagram-symbol identification and whether the source includes this stronger condition cannot be recovered from the extract. Nor does the record state whether its printed target already includes degree completion. The problem must therefore retain its source-defined filtration and target rather than substitute an arbitrary tautological universal map. It explicitly cross-refers to Conjecture 2.17; no equivalence with that extension question is asserted here.
\subsection{Short English statement}
Construct a universal finite-type invariant of virtual knots in the specified directed-diagram space on an interval. Universality should encode the chosen finite-type invariants through finite-degree truncations, but the source's virtual-knot filtration, diagram relations, completion and long-versus-closed convention are missing from this record.
\subsection{Sources}
\begin{itemize}
\item Local UnsolvedMath registry, permanent ID \(10400059\), exact source code \texttt{catalogue entry}, Problem 3.9, attributed to M. Polyak. The supplied title is truncated after ``(M.'' The statement cross-refers to Conjecture 2.17.
\item Ohtsuki, \emph{Problems on invariants of knots and 3-manifolds} (2002), Problem 3.9, PDF page 55: \url{https://msp.org/gtm/2002/04/gtm-2002-04-024s.pdf}. Bold printed headings distinguish 231 numbered items from surrounding references.
\item Extraction pdf-font-structure, accessed 2026-07-29. Source evidence NEEDS\_REVIEW presents an open problem and requires a current primary-literature audit. Supplied OPEN-TRIAGE summary calls this the universal finite-type invariant of virtual knots and says it is not obviously resolved. No further authoritative citation defining the target or resolving the problem is supplied.
\item Dataset attribution: Ulam AI / UnsolvedMath Contributors, CC BY 4.0. Source rights NEEDS\_REVIEW: publicly accessible author/publisher copy, redistribution terms requiring review.
\end{itemize}
\end{document}
